% ===================================================================
% GERMAN SUMMARY
% ===================================================================

\begin{otherlanguage}{german}
	\chapter{Zusammenfassung}
	\normalsize 
	Magnetresonanztomographie (MRT) ist zu einem unverzichtbaren diagnostischen Werkzeug in der Klinik und zu einem der vielseitigsten und leistungsstärksten Bildgebenden Verfahren der Forschung geworden. Auch mehr als fünf Jahrzehnte nach ihrer Einführung erfährt die MRT-Forschung weiterhin erhebliche und kontinuierliche Fortschritte in Technologie, Methodik und Anwendungen. MRT kann eine Vielzahl von Kernen im menschlichen Körper erfassen, von dem reichlich vorhandenen Wasserstoff (\textsuperscript{1}H) bis zu weniger häufig vorkommenden sogenannten „X-Kernen“ wie Deuterium (\textsuperscript{2}H), Kohlenstoff (\textsuperscript{13}C), Natrium (\textsuperscript{23}Na) und Phosphor (\textsuperscript{31}P). Diese Eigenschaften ermöglicht die Messung von Morphologie, Mikrostrukturen und einer Vielzahl physiologischer Parameter, darunter Blutsauerstoffgehalt, Metabolismus, Perfusion, Gewebesteifigkeit, pH-Wert und Temperatur, mit unterschiedlichem Erfolg. Derzeit konzentriert sich ein Großteil der Forschung auf \gls{UHF} ($\geq 7\,T$) MR-Techniken, um die dringend benötigte Sensitivität zu verbessern und spezifische Kontrastmechanismen bei solch hohen Feldstärken zu untersuchen. Diese Arbeit ist eine Zusammenstellung von zwei Projekten, die sich in erster Linie auf die methodische Entwicklung eines schnellen Bildaufnahmeverfahrens namens \gls{bSSFP} zur Untersuchung der Gehirnfunktion bei \gls{UHF} konzentrieren. \Gls{bSSFP} ist bekannt für eine hohe \gls{SNR}-Effizienz und einen charakteristischen $T_2/T_1$-Kontrast mit Flüssigkeitshervorhebung, was es traditionell bei der kardialen Bildgebung bei niedrigeren Feldstärken beliebt macht. 
	
	Das erste Projekt zielt auf die Verbesserung der Aufnahmeeffizienz von bSSFP-basiertem \gls{fMRI} durch schnelle und effiziente nicht-kartesische Trajektorien in Kombination mit paralleler Bildgebung ab. Ziel ist es, eine verzerrungsfreie \gls{fMRI}-Methode mit verbesserter räumlicher Lokalisierung neuronaler Aktivität für \gls{UHF}-Anwendungen als Alternative zu der populären \gls{GE}-\gls{EPI} Methode zu entwickeln. Die mikro-vaskuläre Sensitivität von \gls{bSSFP}, ähnlich der von Spin-Echo-Sequenzen, soll eine bessere Lokalisierung des \gls{BOLD}-Signals an Orten neuronaler Aktivierung ermöglichen. Darüber hinaus ist die inhärent niedrigere Hochfrequenz-Leistungsabgabe von \gls{bSSFP} im Vergleich zu anderen Spin-Echo-Sequenzen bei hohen Feldstärken vorteilhaft. Allerdings sind Bilder, die bei solch hohen räumlichen und zeitlichen Auflösungen aufgenommen werden, hochgradig anfällig für kleinste Fehler in der Bildgebungshardware. Daher wurden Systembedingte Abweichungen, einschließlich Wirbelstromeffekte und räumliche $B_0$-Inhomogenitäten, in das Rekonstruktionsmodell einbezogen um die Bildqualität zu verbessern. Die Nützlichkeit und Flexibilität der entwickelten Aufnahme- und Rekonstruktionsmethode wurde anhand von Submillimeterprotokollen mit begrenzter Abdeckung sowie 1,2\,mm-isotropen Ganzhirn-Messungen bei visuellen Experimenten für funktionelle Bildergebung demonstriert. Darüber hinaus wurden die isolierten Effekte von \gls{TR} und \gls{TE} auf den \gls{bSSFP}-\gls{BOLD}-Kontrast untersucht, indem relative Signaländerung, Signifikanz der Aktivierung sowie thermische und zeitliche Signalstabilität analysiert wurden.
	
	Im zweiten Projekt wurde die Kombination von \gls{bSSFP} Messungen bei \gls{UHF} untersucht, um die Sensitivität von \gls{DMI} zu verbessern. \Gls{DMI} ist eine neuartige molekulare \gls{MRSI}-Technik, die den Metabolismus mithilfe eines nicht-radioaktiven, deuterium-markierten exogenen Tracers, typischerweise $[6,6^2-^2H_2]$-Glukose, untersucht. Im Gegensatz zum Goldstandard \textsuperscript{18}F-FDG, der in bei \gls{PET} verwendet wird, radioaktiv ist und nicht metabolisiert wird, kann deuterierte Glukose metabolisiert werden und liefert zusätzliche Informationen über nachgeschaltete Stoffwechselwege. Dies ist ein entscheidender Vorteil im Gehirn, wo die allgemeine Glukoseaufnahme typischerweise sehr hoch ist. Trotz dieser Vorteile machen das niedrige gyromagnetische Verhältnis von Deuterium und die geringen Metabolitenkonzentrationen die Technik oft unpraktisch. Positiv ist jedoch, dass das kurze $T_1$, resultierend aus der ausgeprägten quadrupolaren Relaxation der Deuteriumkerne, sowie das schmale Deuteriumspektrum eine schnellere Signal-Mittelung unterstützen. Diese Vorteile, kombiniert mit den \gls{SNR}-Verbesserungen durch \gls{UHF} und \gls{bSSFP} – letzteres resultierend aus dem großen $T_2/T_1$-Verhältnis von Deuterium-Metaboliten – verschaffen \gls{DMI} einen Wettbewerbsvorteil gegenüber anderen MR-Molekularbildgebungsverfahren wie \gls{CEST}, Protonen (\textsuperscript{1}H) \gls{MRSI} und \textsuperscript{13}C MRI. Herausforderungen im Zusammenhang mit großen $B_0$-Inhomogenitäten im menschlichen Kopf wurden mithilfe von Phasenzyklen adressiert, wenn auch auf Kosten des SNR, um off-resonanz-unempfindliche Ganzhirn-Metabolitenkarten zu erhalten. Die in dieser Arbeit eingeführte neuartige Verarbeitungstechnik zur Schätzung von Metabolitenamplituden aus phasen-zyklischen \gls{bSSFP}-Daten bietet größere Flexibilität, da sie ohne explizites Wissen über die räumliche Off-Resonanz-Karte, Flipwinkelkarte oder Relaxationszeiten der Deuterium-Metaboliten funktioniert.
	
	Insgesamt hebt diese Arbeit sowohl das Potenzial als auch die Herausforderungen neuartiger \gls{UHF}-Steady-State-Techniken für die Weiterentwicklung funktioneller und metabolischer Untersuchungen des menschlichen Gehirns hervor.
	
\end{otherlanguage}


%%% Local Variables:
%%% mode: latex
%%% TeX-master: "../thesis"
%%% End:
